% !TEX root = ../main.tex
% ===========================================================================
% SOURCES: software/export/program.txt (the 64-instruction listing)
%          scripts/gen_program.py (the encoder - the authority on the format)
%          hardware/rtl/control/top_seq.v (the decoder)
%
% Generate the encoding table from gen_program.py rather than transcribing it.
% An appendix that disagrees with the assembler is worse than no appendix.
% ===========================================================================
\chapter*{Appendix A: Accelerator Instruction Set}
\label{appendix:isa}
\addcontentsline{toc}{chapter}{Appendix A: Accelerator Instruction Set}

\section{Instruction Format}
\label{sec:isa-encoding}
% The bit layout: opcode, operands, field widths, and which fields each opcode
% actually reads.

\section{Opcodes}
\label{sec:isa-opcodes}
% The six opcodes, one subsection or table row each: what it computes, which
% engine it drives, which feeder, and its latency contract
% (Section~\ref{sec:latency-contract}).
% Note that `save' costs zero cycles and zero instructions, which is why the
% manifest's \thManifestOps{} ops compile to \thInstructions{} instructions.

\section{The Compiled Program}
\label{sec:isa-program}
% The full \thInstructions-instruction listing for this network, from
% software/export/program.txt. This is the right thing to put in an appendix:
% it is generated, it is short, and it lets a reader follow
% Section~\ref{sec:sequencer} against real code.
